\pagebreak

\section{Xilinx GTP library - xilinx/gtp\_dual.3pl}


    \index{library!xilinx/gtp\_dual.3pl}
    This library creates GTP dual tiles.
    
    This library must be incorporated by using a {\bf module}
    statement -
    \begin{lstlisting}[gobble=4, language=threepl]
       module "xilinx/gtp_dual.3pl"
    \end{lstlisting}
            
    The struct {\bf gtp\_tile\_params} describes the pins and some control
    parameters for a GTP tile (GTP channel pair).

    \begin{lstlisting}[gobble=4, language=threepl]
       struct(global.gtp_tile_params,
           loc,                                        "str",
           {refclkp_pin, refclkn_pin},                 "str",
           {boost0, boost1},                           "log",
           {txdiffctrl0, txdiffctrl1},                 "int",
           {txpreemphasis0, txpreemphasis1},           "int"
       );
    \end{lstlisting}

    The following structs are created for GTP data input and output -

    \begin{lstlisting}[gobble=4, language=threepl]
       struct(global.gt_16,
           data,   "bits:16",  // data
           k,      "log",      // true if top 8 bits of data form a K-character
           contin, "log"       // this is a continuation word
       );
    \end{lstlisting}

    \begin{lstlisting}[gobble=4, language=threepl]
       struct(global.gt_16_e,
           data,   "bits:16"   // data
           k,      "log",      // true if top 8 bits of data form a K-character
           first,  "log",      // previous word was an idle
           error,  "log"       // disparity error or not-in-table error
       );
    \end{lstlisting}

    \begin{lstlisting}[gobble=4, language=threepl]
       struct(global.gt_16_ee,
           data,       "bits:16",  // data
           k,          "log",      // true if top 8 bits of data form a K-character
           first,      "log",      // previous word was an idlgtp\_duale
           disperr,    "[2]log",   // disparity errors
           niterr,     "[2]log"    // not-in-table errors
       );
    \end{lstlisting}
    
    The global structure {\bf ds16} is also used and it is created in
    {\bf stdlib.3pl}. It is 
    \begin{lstlisting}[gobble=4, language=threepl]
       struct(global.ds16,
           data,   "bits:16",  // data
           k,      "log"       // true if top 8 bits of data form a K-character
       );
    \end{lstlisting}


    \subsection{library modules}

        \subsubsection{gtp\_dual\_16 - a 16-bit dual high-speed serial I/O channel tile}\label{sec:lmgtpdual16}

            {\bf Library: xilinx/gtp\_dual.3pl}

            {\bf Prototype:} \\
            \vsp
            \begin{lstlisting}[gobble=12, language=threepl]
               global module
               gtp_dual_16 (
                   queue(out0),
                   queue(out1),
                   clock(usrclk2),
                   value({rxdiag0, rxdiag1}, "[10]log"),
                   value({txdiag0, txdiag1}, "[2]log")
                   <-
                   var(tile, gtp_tile_params),
                   float(freq, 156.25),
                   int(run_length),
                   queue(in0),
                   queue(in1),
                   value({txpowerdown0, txpowerdown1, rxpowerdown0, rxpowerdown1}, "uint:2", 0),
                   value({gtpreset, rxreset0, rxreset1}, "log", false),
                   value({loopback0, loopback1}, "uint:3", 0),
                   value({diff0, diff1, pe0, pe1}, "uint:3", 0),
                   log(los_fsm)
               ) {}
            \end{lstlisting}

            {\bf Output parameters:} \\
            {\bf out0} is the data output queue for GTP0.\\
            {\bf out1} is the data output queue for GTP1.\\
            {\bf usrclk2} is the user interface clock.\\
            {\bf rxdiag0} is an array of diagnostic signals for the GTP0 receiver.\\
            {\bf rxdiag1} is an array of diagnostic signals for the GTP1 receiver.
            {\bf txdiag0} is an array of diagnostic signals for the GTP0 transmitter.\\
            {\bf txdiag1} is an array of diagnostic signals for the GTP1 transmitter.
            
            {\bf Input parameters:} \\
            {\bf tile} indicates which GTP\_DUAL tile to use.\\
            {\bf run\_length} is the maximum data run length.\\
            {\bf in0} is the data input queue fro GTP0.\\
            {\bf in1} is the data input queue fro GTP1.\\
            {\bf txpowerdown0} control GTP0 transmitter powerdown.\\
            {\bf txpowerdown1} control GTP1 transmitter powerdown.\\
            {\bf rxpowerdown0} control GTP0 receiver powerdown.\\
            {\bf rxpowerdown1} control GTP1 receiver powerdown.\\
            {\bf gtpreset} when true resets the entire GTP\_DUAL tile.\\
            {\bf rxreset0} when true resets the GTP0 receiver.\\
            {\bf rxreset1} when true resets the GTP1 receiver.\\
            {\bf loopback0} controls the GTP0 loopback mode.\\
            {\bf loopback1} controls the GTP1 loopback mode.\\
            {\bf diff0} overrides the GTP0 differential output voltage table entry.\\
            {\bf diff1} overrides the GTP1 differential output voltage table entry.\\
            {\bf pe0} overrides the GTP0 transmitter pre-emphasis table entry.\\
            {\bf pe1} overrides the GTP1 transmitter pre-emphasis table entry.\\
            {\bf los\_fsm} when true instantiates the loss-of-sync FSM

            {\bf In-line target execution:} \\
            no

            {\bf Description:} \\
            This module is both an input and an output interface to a
            GTP\_dual GTP tile. The data width is 16 bit. A custom
            protocol is used. This calls {\bf
            gtp\_dual()} \autorefph{sec:lpgtpdual}, described below.
            The default external clock frequency is 156.25 MHz (serial
            rate 3.125Gbit/s) but lower clock frequencies can be used.
            
            The particular GTP\_DUAL tile is selected by the argument to input
            parameter {\bf tile}. This parameter has the following type -
            
            \begin{lstlisting}[gobble=12, language=threepl]
               struct(global.gtp_tile_params,
                   loc,                                        "str",
                   {refclkp_pin, refclkn_pin},                 "str",
                   {boost0, boost1},                           "log",
                   {txdiffctrl0, txdiffctrl1},                 "int",
                   {txpreemphasis0, txpreemphasis1},           "int"
               );
            \end{lstlisting}

            where the fields are -
            \blist
            \item   loc - the tile location, e.g. "GTP\_DUAL\_X0Y4"
            \item   refclkp\_pin - the external reference clock positive pin
            \item   refclkn\_pin - the external reference clock negative pin
            \item   boost0 - GTP0 TX pre-emphasis boost
            \item   boost1 - GTP1 TX pre-emphasis boost
            \item   txdiffctrl0 - GTP0 TX differential Voltage
            \item   txdiffctrl1 - GTP1 TX differential Voltage
            \item   txpreemphasis0 - GTP0 TX pre-emphasis
            \item   txpreemphasis1 - GTP1 TX pre-emphasis
            \blend
            
            {\bf boost0} and {\bf boost1} are logical values which select one of
            two sets of levels of transmitter high frequency pre-emphasis (see
            {\bf txpreemphasis0} and {\bf txpreemphasis1} description below)

            {\bf txdiffctrl0} and {\bf txdiffctrl1} are 3 bit values which
            control the transmitter output voltage swing. These table values can
            be overridden by input parameters {\bf diff0} and {\bf diff1}. The
            values are -
            
            \begin{tabular}{| c | c |}
            \hline
            3-bit value      & transmitter differential swing (mV)  \\
            \hline
            \hline
            000 & 1100 \\
            \hline
            001 & 1050 \\
            \hline
            010 & 1000 \\
            \hline
            011 & 900 \\
            \hline
            100 & 800 \\
            \hline
            101 & 600 \\
            \hline
            110 & 400 \\
            \hline
            111 & 0 \\
            \hline
            \end{tabular}

            {\bf txpreemphasis0} and {\bf txpreemphasis1} are 3 bit values which
            control the transmitter high frequency pre-emphasis. These table
            values can be overridden by input parameters {\bf pe0} and {\bf
            pe1}. The values are -
            
            \begin{tabular}{| c | c | c |}
            \hline
            3-bit value & \multicolumn{2}{c|}{transmitter pre-emphasis (\%)} \\
                        & pre-emphasis boost off & pre-emphasis boost on \\
                        & boost = false & boost = true \\
            \hline
            \hline
            000 & 2 & 3 \\
            \hline
            001 & 2 & 3 \\
            \hline
            010 & 2.5 & 4 \\
            \hline
            011 & 4.5 & 10.5 \\
            \hline
            100 & 9.5 & 18.5 \\
            \hline
            101 & 16 & 28 \\
            \hline
            110 & 23 & 39 \\
            \hline
            111 & 31 & 52 \\
            \hline
            \end{tabular}
            
            Input parameter {\bf run\_length} is an immediate integer which
            specifies the maximum contiguous data run for the transmitter. If
            data input is continuously available such that this length would be
            exceeded (and the data {\bf contin} field is not true - see above)
            then an idle sequence is inserted in the output stream. When no data
            are available for transmission a continuous stream of idles is
            transmitted. These idles are used by the receiver to lock initially
            and to continue checking that lock is maintained. If receiver lock
            is lost then lock can only be regained at the next idle. In the
            custom protocol used in this interface the idle sequence is a single
            16-bit word containing a comma byte. If no argument is supplied
            to parameter {\bf run\_length} then idles will only be inserted
            when the input queue is unavailable.

            Input parameters {\bf in0} and {\bf in1} are the input data queues
            for the transmitters. If a transmitter is not used the associated
            argument is omitted. The type is either the struct {\bf ds16}
            (defined in stdlib.3pl) or the struct {\bf gt\_16}. The most
            significant byte of the word may be a K-character, which is
            signified by setting field {\bf k} true. The byte must be one of the
            allowed K-characters (this is not checked!). For type {\bf gt\_16}
            the {\bf contin} field should be true if this word is a continuation
            word, i.e. it is not the first word in a block. This ensures that
            despite the run length constraint a block will not be split by
            inserted idles. At the receiving end the first word following an
            idle can reset a block boundary.
            
            Input parameters {\bf txpowerdown0}, {\bf txpowerdown1}, {\bf
            rxpowerdown0} and {\bf rxpowerdown1} control transmitter and
            receiver powerdown - see Xilinx ''Virtex-5 FPGA RocketIO GTP Tranceiver -
            User Guide`` UG196.
            
            Input parameter {\bf gtpreset} when high resets the entire GTP\_DUAL
            tile including the PLL.
            
            Input parameters {\bf rxreset0} and {\bf rxreset1} when high reset
            the associated receiver.
            
            Input parameters {\bf loopback0} and {\bf loopback1} control the
            loopback mode of each GTP - see Xilinx ''Virtex-5 FPGA RocketIO GTP
            Tranceiver - User Guide`` UG196.

            Output parameters {\bf out0} and {\bf out0} are the output data
            queues for the receivers. If the receiver is not used the argument is
            omitted. The type can be struct {\bf ds16}, struct {\bf gt\_16\_e}
            or  struct {\bf gt\_16\_ee} depending on the level of error detail
            required. If a receiver loses sync or byte alignment there is no
            data output. A word in which a disparity error or a not-in-table
            error occurs is written to the output queue but the error field is
            true. If the top byte of a word is a K-character the {\bf k} field
            is true. If the {\bf first} field is true then this word follows an
            idle, which information may be used to reset a block boundary. {\bf
            Important note - the output queues must not block! The destination
            module must read them at a rate that precludes them ever becoming
            full, if necessary reading them to null. Generally the read clock
            domain must have an equal or higher frequency than the user
            interface clock {\bf ruclk}, but this may be relaxed if it is known
            that the serial data source has a guaranteed maximum data rate that
            would allow a lower queue read clock domain.}

            Output parameter {\bf usrclk2} is the user interface clock (default
            156.25MHz).
            
            Output parameters {\bf rxdiag0} and {\bf rxdiag1} provide ten
            diagnostic outputs for the associated receiver. These are -
            \blist
            \item   diag[0] - receiver synchronisation FSM bit 1, if los\_fsm true
            \item   diag[1] - receiver synchronisation FSM bit 0, if los\_fsm true
            \item   diag[2] - is true when the upper byte of the current output word
                    has a disparity error
            \item   diag[3] - is true when the lower byte of the current output word
                    has a disparity error
            \item   diag[4] - is true when the upper byte of the current output word
                    is a comma
            \item   diag[5] - is true when the lower byte of the current output word
                    is a comma
            \item   diag[6] - rxbytisaligned, which is true when the idles are
                    word boundary aligned
            \item   diag[7] - rxcommadet, which is true when a positive or negative comma
                    is being received
            \item   diag[8] - ok, which is true when synchronisation is locked and the
                    byte pair is aligned properly to a 16 bit word
            \item   diag[9] - idle, which is true when the receiver is idle
            \blend
            {\bf diag[0}] and {\bf diag[1]} are both false if {\bf los\_fsm} is
            false since the loss-of-sync FSM will not have been instantiated.
            These outputs might be connected to status LEDs to indicate that the
            receiver is locked and functioning correctly. diag[5] will be true
            while the receiver is receiving idles but will be false when data
            are received.

            Output parameters {\bf txdiag0} and {\bf txdiag1} provide two
            diagnostic outputs for the associated transmitter. These are -
            \blist
            \item   diag[0] - is\_k, which is true when a data word with the
                    K bit set is being transmitted
            \item   diag[1] - idle, which is true when the receiver is idle
            \blend




    \subsection{library procedures}
    

        \subsubsection{gtp\_dual - create a GTP dual high-speed serial I/O channel tile}\label{sec:lpgtpdual}

            {\bf Library: xilinx/gtp\_dual.3pl}

            {\bf Prototype:} \\
            \vsp
            \begin{lstlisting}[gobble=12, language=threepl]
               global procedure
               gtp_dual (
                   // outputs
                   clock({refclkout, rxrecclk0, txoutclk0, rxrecclk1, txoutclk1}),
                   value({
                       drdy, phystatus0, phystatus1, plllkdet, resetdone0, resetdone1,
                       rxbyteisaligned0, rxbyteisaligned1, rxbyterealign0, rxbyterealign1,
                       rxchanbondseq0, rxchanbondseq1, rxchanisaligned0, rxchanisaligned1,
                       rxchanrealign0, rxchanrealign1, rxcommadet0, rxcommadet1, rxelecidle0,
                       rxelecidle1, rxoversampleerr0, rxoversampleerr1, rxprbserr0, rxprbserr1,
                       rxvalid0, rxvalid1},        
                       "log"
                   ),
                   value(rxconfigout, "bits:16"),  // port do
                   value({rxbufstatus0, rxbufstatus1}, "uint:3"),
                   value({rxchbondo0, rxchbondo1}, "[3]log"),
                   value({rxclkcorcnt0, rxclkcorcnt1}, "uint:3"),
                   value({rxdata0, rxdata1}, "[2]bits:8"),
                   value({rxdisperr0, rxdisperr1, rxnotintable0, rxnotintable1,
                          rxlossofsync0, rxlossofsync1, rxrundisp0, rxrundisp1,
                          rxchariscomma0, rxchariscomma1, rxcharisk0, rxcharisk1,
                          txkerr0, txkerr1, txrundisp0, txrundisp1},
                          "[2]log"),
                   value({rxstatus0, rxstatus1}, "uint:3"),
                   value({txbufstatus0, txbufstatus1}, "[2]log")

                   <-

                   // attributes
                   str({loc, refclkp_pin, refclkn_pin}),
                   log({
                       ac_cap_dis_0, ac_cap_dis_1, chan_bond_seq_2_use_0,
                       chan_bond_seq_2_use_1, clk_cor_insert_idle_flag_0,
                       clk_cor_insert_idle_flag_1, clk_cor_keep_idle_0, clk_cor_keep_idle_1,
                       clk_cor_precedence_0, clk_cor_precedence_1, clk_cor_seq_2_use_0,
                       clk_cor_seq_2_use_1, clk_correct_use_0, clk_correct_use_1,
                       comma_double_0, comma_double_1, dec_mcomma_detect_0,
                       dec_mcomma_detect_1, dec_pcomma_detect_0, dec_pcomma_detect_1,
                       dec_valid_comma_only_0, dec_valid_comma_only_1, mcomma_detect_0,
                       mcomma_detect_1, oversample_mode, pci_express_mode_0,
                       pci_express_mode_1, pcomma_detect_0, pcomma_detect_1, rcv_term_gnd_0,
                       rcv_term_gnd_1, rcv_term_mid_0, rcv_term_mid_1, rcv_term_vttrx_0,
                       rcv_term_vttrx_1, rx_buffer_use_0, rx_buffer_use_1,
                       rx_decode_seq_match_0, rx_decode_seq_match_1, rx_loss_of_sync_fsm_0,
                       rx_loss_of_sync_fsm_1, termination_ovrd, tx_buffer_use_0,
                       tx_buffer_use_1
                   }),
                   log(tx_diff_boost_0, true),
                   log(tx_diff_boost_1, true),
                   int({
                       align_comma_word_0, align_comma_word_1, chan_bond_1_max_skew_0,
                       chan_bond_1_max_skew_1, chan_bond_2_max_skew_0, chan_bond_2_max_skew_1,
                       chan_bond_level_0, chan_bond_level_1, chan_bond_seq_1_1_0,
                       chan_bond_seq_1_1_1, chan_bond_seq_1_2_0, chan_bond_seq_1_2_1,
                       chan_bond_seq_1_3_0, chan_bond_seq_1_3_1, chan_bond_seq_1_4_0,
                       chan_bond_seq_1_4_1, chan_bond_seq_1_enable_0, chan_bond_seq_1_enable_1,
                       chan_bond_seq_2_1_0, chan_bond_seq_2_1_1, chan_bond_seq_2_2_0,
                       chan_bond_seq_2_2_1, chan_bond_seq_2_3_0, chan_bond_seq_2_3_1,
                       chan_bond_seq_2_4_0, chan_bond_seq_2_4_1, chan_bond_seq_2_enable_0,
                       chan_bond_seq_2_enable_1, clk_cor_seq_1_1_0, clk_cor_seq_1_1_1,
                       clk_cor_seq_1_2_0, clk_cor_seq_1_2_1, clk_cor_seq_1_3_0,
                       clk_cor_seq_1_3_1, clk_cor_seq_1_4_0, clk_cor_seq_1_4_1,
                       clk_cor_seq_1_enable_0, clk_cor_seq_1_enable_1, clk_cor_seq_2_1_0,
                       clk_cor_seq_2_1_1, clk_cor_seq_2_2_0, clk_cor_seq_2_2_1,
                       clk_cor_seq_2_3_0, clk_cor_seq_2_3_1, clk_cor_seq_2_4_0,
                       clk_cor_seq_2_4_1, clk_cor_seq_2_enable_0, clk_cor_seq_2_enable_1,
                       com_burst_val_0, com_burst_val_1, comma_10b_enable_0, comma_10b_enable_1,
                       mcomma_10b_value_0, mcomma_10b_value_1, oobdetect_threshold_0,
                       oobdetect_threshold_1, pcomma_10b_value_0, pcomma_10b_value_1,
                       sata_burst_val_0,
                       sata_burst_val_1, sata_idle_val_0, sata_idle_val_1, termination_ctrl
                   }),
                   int({   pcs_com_cfg, pma_cdr_scan_0, pma_cdr_scan_1, pma_rx_cfg_0,
                           pma_rx_cfg_1, prbs_err_threshold_0, prbs_err_threshold_1, 
                           trans_time_from_p2_0, trans_time_from_p2_1, trans_time_non_p2_0,
                           trans_time_non_p2_1, trans_time_to_p2_0, trans_time_to_p2_1}),
                   int(chan_bond_seq_len_0,    1),
                   int(chan_bond_seq_len_1,    1),
                   int(clk25_divider,          1),
                   int(clk_cor_adj_len_0,      1),
                   int(clk_cor_adj_len_1,      1),
                   int(clk_cor_det_len_0,      1),
                   int(clk_cor_det_len_1,      1),
                   int(clk_cor_max_lat_0,      3),
                   int(clk_cor_max_lat_1,      3),
                   int(clk_cor_min_lat_0,      3),
                   int(clk_cor_min_lat_1,      3),
                   int(clk_cor_repeat_wait_0,  0),
                   int(clk_cor_repeat_wait_1,  0),
                   int(oob_clk_divider,        1),
                   int(pll_divsel_fb,          1),
                   int(pll_divsel_ref,         1),
                   int(pll_rxdivsel_out_0,     1),
                   int(pll_rxdivsel_out_1,     1),
                   int(pll_txdivsel_comm_out,  1),
                   int(pll_txdivsel_out_0,     1),
                   int(pll_txdivsel_out_1,     1),
                   int(rx_los_invalid_incr_0,  1),
                   int(rx_los_invalid_incr_1,  1),
                   int(rx_los_threshold_0,     4),
                   int(rx_los_threshold_1,     4),
                   int(sata_max_burst_0,       7),
                   int(sata_max_burst_1,       7),
                   int(sata_max_init_0,       22),
                   int(sata_max_init_1,       22),
                   int(sata_max_wake_0,        7),
                   int(sata_max_wake_1,        7),
                   int(sata_min_burst_0,       5),
                   int(sata_min_burst_1,       5),
                   int(sata_min_init_0,       12),
                   int(sata_min_init_1,       12),
                   int(sata_min_wake_0,        4),
                   int(sata_min_wake_1,        14),
                   str({   chan_bond_mode_0, chan_bond_mode_1, rx_slide_mode_0,
                           rx_slide_mode_1, rx_status_fmt_0, rx_status_fmt_1, rx_xclk_sel_0,
                           rx_xclk_sel_1   }),
                   // str(pma_com_cfg),
                   // inputs
                   clock({
                       txusrclk0, txusrclk20, rxusrclk0, rxusrclk20,
                       txusrclk1, txusrclk21, rxusrclk1, rxusrclk21,
                       clkin, dclk
                   }),
                   value({den, dwe, gtpreset, intdatawidth, plllkdeten, pllpowerdown,
                          prbscntreset0, prbscntreset1, refclkpwrdnb, rxbufreset0, rxbufreset1,
                          rxcdrreset0, rxcdrreset1, rxcommadetuse0, rxcommadetuse1,
                          rxdatawidth0, rxdatawidth1, rxdec8b10buse0, rxdec8b10buse1,
                          rxelecidlereset0, rxelecidlereset1, rxenchansync0, rxenchansync1,
                          rxenelecidleresetb, rxeneqb0, rxeneqb1, rxenmcommaalign0,
                          rxenmcommaalign1, rxenpcommaalign0, rxenpcommaalign1,
                          rxensamplealign0, rxensamplealign1, rxpmasetphase0, rxpmasetphase1,
                          rxpolarity0, rxpolarity1, rxreset0, rxreset1, rxslide0, rxslide1,
                          txcomstart0, txcomstart1, txcomtype0, txcomtype1, txdatawidth0,
                          txdatawidth1, txdetectrx0, txdetectrx1, txelecidle0, txelecidle1,
                          txenc8b10buse0, txenc8b10buse1, txenpmaphasealign, txinhibit0,
                          txinhibit1, txpmasetphase, txpolarity0, txpolarity1, txreset0, txreset1
                          }, "log"
                   ),
                   value(daddr, "uint:7"),
                   value(txconfigin, "bits:16"),   // port di
                   value({rxchbondi0, rxchbondi1}, "[3]log"),
                   value({rxeqpole0, rxeqpole1}, "uint:4"),
                   value({txbypass8b10b0, txbypass8b10b1, txchardispmode0, txchardispmode1, txchardispval0, txchardispval1,
                          txcharisk0, txcharisk1},
                          "[2]log"),
                   value({rxeqmix0, rxeqmix1}, "uint:2"),
                   value({txdata0, txdata1}, "[2]bits:8"),
                   value({txenprbstst0, rxenprbstst0,
                          txenprbstst1, rxenprbstst1},
                         "uint:2"),  
                   value({txpowerdown0, rxpowerdown0,
                          txpowerdown1, rxpowerdown1},
                         "uint:2"),
                   value({txpreemphasis0, txpreemphasis1, txdiffctrl0, txdiffctrl1,
                          txbufdiffctrl0, txbufdiffctrl1, loopback0, loopback1},
                          "uint:3")
               ) {
            \end{lstlisting}

            {\bf Output parameters:} \\
            See Xilinx documentation.

            {\bf Input parameters:} \\
            See Xilinx documentation.

            {\bf In-line target execution:} \\
            no

            {\bf description:} \\
            This procedure creates a dual GTP tile. For a simpler interface
            see {\bf gtp\_dual\_16} \autorefph{sec:lmgtpdual16} above.




    \subsection{library functions}
    
        None.






    \subsection{library classes}
    
        None.
